\section{Exact four-wave reduction and audit of normal powers}
\label{sec:constraint-wave-reduction}

This is a coefficient-level audit of the normal structure used in
Section~\ref{sec:constraint-normal-bootstrap}. Put
\[
 l=1/a,\quad d_0=(aS)^{-1},\quad
 d=\kappa_1(2+\kappa_2),\quad D=3\kappa_1(1+\kappa_2),
 \quad c=2l-\kappa_1,\quad C=A+\Omega,
 \quad D_0=\partial_t-B\cdot\partial,
\]
and write $\mathcal P=D_0^2-A^2\Delta$,
$s_z=\partial_i z_i$, $g_z=G\cdot z$. All vector components in the following
two equations are Cartesian. The exact corrected subsidiary equations give
\begin{align}
H&=(2\Omega/A)D_0T-2\Omega^2s_z+2dT,\\
M_i&=(\Omega/A)(D_0z_i+2lz_i)-\partial_iT-cz_i.
\end{align}
Substitution into the exact Hamiltonian and momentum equations, using
$[D_0,\partial_i]=-l\partial_i$, gives the scalar equation
\begin{align}
0={}&\mathcal P T+
 \frac{wC+A(d+6l)}\Omega D_0T
 -(2A+\Omega)G\cdot D_0z\\
&+[-2A^2w+3lA\Omega-A^2(c+2l)]s_z
 +\frac{A(3A+2\Omega)}\Omega G\cdot\partial T\\
&+\left[-2l(2A+\Omega)+\frac{Ac(3A+2\Omega)}\Omega\right]g_z
 -\frac{2lA^2(D-3d)}{\Omega^2}T .
\label{eq:constraint-scalar-wave-exact}
\end{align}
For the vector equation define the scalar coefficient
\[
 Z_c=\frac{2lwC}{A}+2l^2(\Omega/A+3)
 -lc(1+3A/\Omega)+2G^2+\frac{20}{3}Ad_0+
 \frac{6AG^2}\Omega.
\]
Then the exact vector equations are
\begin{align}
0={}&\mathcal P z_i+
 \left[\frac{wC}\Omega+3l+\frac{A(3l-c)}\Omega\right]D_0z_i
 -\frac{CG_i}{\Omega^2}D_0T-\frac{2A^2G_i}\Omega s_z\\
&+A G\cdot\partial z_i+
 \frac{A(10A/3+\Omega)}\Omega G_j\partial_i z_j
 +\frac A\Omega Z_c z_i\\
&+\left[\frac{2A}{3\Omega}+\frac{2A^2}{\Omega^2}\right]G_i g_z
 -\frac{A^2(3l-d+2D/3)}{\Omega^2}\partial_iT
 -\frac{2\kappa_1AC}{\Omega^3}G_iT .
\label{eq:constraint-vector-wave-exact}
\end{align}
The $\Omega^{-3}T$ term is real; it must not be omitted. It is balanced by
the different powers in $T=\Omega^2\tau$ and $z=\Omega\zeta$.

One derivation of the vector identity first uses
\begin{align}
 D_p^jQ^p_{ij}-D_i\operatorname{tr}_pQ^p
 ={}&A\Omega(\Delta z_i-\partial_i s_z)
 -\Omega G\cdot\partial z_i
 -(10A/3+\Omega)G_j\partial_i z_j+2CG_is_z\\
&-\left[2G^2+\tfrac{20}{3}Ad_0+
                  \tfrac{6AG^2}\Omega\right]z_i
 -\left[\tfrac23+\tfrac{2A}\Omega\right]G_i g_z\\
&+\frac{2DA}{3\Omega}\partial_iT
 -\frac{2DC}{3\Omega^2}G_iT .
\end{align}
The mixed terms $D_0\partial_iT$ and $\partial_iD_0T$ cancel, as do
$\partial_i s_z$ from the Hamiltonian derivative and the tensor divergence.
Likewise $D_0s_z$ cancels $\partial_iD_0z_i$ in the scalar equation. Thus no
second-order cross principal operator survives.
The commutator in $D_0M_i$ also contributes $+l\partial_iT$ on the left;
it cancels the corresponding momentum-reaction contribution. There is
therefore no extra $-Al\partial_iT/\Omega$ in
\eqref{eq:constraint-vector-wave-exact}. A lightweight independent
automatic-differentiation evaluation of the nested and expanded identities
at a generic nonradial point detected and corrected precisely this term.

\paragraph{Audit of the normal eight-system.}
In spherical coordinates, with Cartesian dependent components,
\[
 \mathcal P=\partial_t^2-2B_r\partial_t\partial_r
 -\Omega^2\partial_r^2+
 (B_r\partial_rB_r-2A^2/r)\partial_r
 -A^2r^{-2}\Delta_{\mathbb S^2}.
\]
Since $\partial_r=-\gamma\partial_\Omega$,
$\gamma=r/(aS)$, substituting the balanced powers and multiplying the
balanced equations by $\Omega^2$ has exactly the following consequences:
\begin{itemize}
\item The highest radial term is $-\gamma^2\Omega^3\partial_\Omega N$,
where $N=\Omega\partial_\Omega(\tau,\zeta)$. The term obtained by
differentiating $1/\Omega$ is a smooth $O(\Omega^2)N$ coefficient.
\item The only derivative of $N$ in time/angular directions is
$2B_r\gamma\Omega\partial_tN$. There is no angular derivative of $N$.
\item All other first time/angular derivatives of $(\tau,\zeta)$ carry
at least one power of $\Omega$. Second time/angular derivatives carry
$\Omega^2$. Every zero-order and $N$ coefficient is smooth up to scri.
\item Passing to normal/tangential vector components changes only smooth
angular lower-order coefficients: the frame is constant along each radial
ray. It does not introduce a mixed normal/angular principal term.
\end{itemize}
These statements follow term by term from
\eqref{eq:constraint-scalar-wave-exact}--\eqref{eq:constraint-vector-wave-exact};
there is no residual $\Omega^{-1}$ in a derivative coefficient of the
eight-system after the displayed multiplication.

At scri, in $a=1$ units and for $\kappa_1=4$, $\kappa_2=-99/100$, the
scalar/normal $N$ and value matrices in
$-\gamma^2\Omega^3N_\Omega+LN+Kv+\cdots=0$ are
\[
 L_0=A_{\scri}^2
 \begin{pmatrix}-151/25&-4\\-49/25&-8/3\end{pmatrix},\qquad
 K_0=A_{\scri}^2
 \begin{pmatrix}298/25&2\\102/25&34/3\end{pmatrix}.
\]
Each transverse block is $L_0=-4A_{\scri}^2$,
$K_0=8A_{\scri}^2$. Hence $-L_0$ has positive symmetric part and
\[
 -L_0^{-1}K_0=
 \begin{pmatrix}58/31&-150/31\\24/155&242/31\end{pmatrix},
\]
with transverse entries $2$. Equivalently $L_0=2A_{\scri}^2lJ_{gg}$ with
the appropriately dimensioned glancing matrix. The fast coercivity therefore
uses the same wave metric as the previously certified energy, without an
unverified change of sign or target metric.

Finally, setting $F=-L^{-1}K$ and $R_g=N-Fv$ introduces only
$\Omega^3F_\Omega v$ and $\Omega^2F(R_g+Fv)$ from the highest normal term.
The first is higher order; the $R_g$ term can be absorbed into a smooth
coercive fast zero-order matrix in a smaller collar. The mixed time term
gives $\Omega\partial_tR_g$ and $\Omega F\partial_tv$. Thus the fast forcing
order/power count required by the finite graph bootstrap is verified by the
exact equations, not merely assumed from the indicial roots.

